\chapter{Activation Functions and Normalization}
\label{chap:activation_normalization}

\begin{tldr}
Activation functions introduce nonlinearity; normalization techniques stabilize training. This chapter covers the evolution of activations (sigmoid → ReLU → GELU → SwiGLU), normalization methods (BatchNorm, LayerNorm, RMSNorm), and when to use each. Understanding why these components work—and their failure modes—is essential for debugging training issues.
\end{tldr}

\section{Activation Functions}
\label{sec:activations}

\subsection{Why Nonlinearity?}

Without activation functions, deep networks collapse to linear models:
\begin{equation}
f(\vx) = \mW_L \cdots \mW_2 \mW_1 \vx = \mW_{eff} \vx
\end{equation}

Nonlinearity enables learning complex decision boundaries and representations.

\subsection{Classical Activations}

\subsubsection{Sigmoid}

\begin{mathresult}{Sigmoid}
\begin{equation}
\sigma(x) = \frac{1}{1 + e^{-x}}, \quad \sigma'(x) = \sigma(x)(1 - \sigma(x))
\end{equation}
\end{mathresult}

\textbf{Properties}:
\begin{itemize}
    \item Output range: $(0, 1)$
    \item Smooth, differentiable everywhere
    \item Maximum gradient: 0.25 (at $x=0$)
\end{itemize}

\textbf{Problems}:
\begin{itemize}
    \item \textbf{Vanishing gradients}: For $|x| > 4$, gradient $\approx 0$
    \item \textbf{Not zero-centered}: Outputs always positive, causes zig-zag gradients
    \item \textbf{Expensive}: Exponential computation
\end{itemize}

\subsubsection{Tanh}

\begin{equation}
\tanh(x) = \frac{e^x - e^{-x}}{e^x + e^{-x}} = 2\sigma(2x) - 1
\end{equation}

\textbf{Improvement over sigmoid}: Zero-centered output $(-1, 1)$

\textbf{Still has}: Vanishing gradients for $|x| > 2$

\subsection{ReLU Family}

\subsubsection{ReLU (Rectified Linear Unit)}

\begin{mathresult}{ReLU}
\begin{equation}
\text{ReLU}(x) = \max(0, x), \quad \text{ReLU}'(x) = \begin{cases} 1 & x > 0 \\ 0 & x \leq 0 \end{cases}
\end{equation}
\end{mathresult}

\textbf{Advantages}:
\begin{itemize}
    \item No vanishing gradient for $x > 0$
    \item Sparse activations (biological plausibility)
    \item Computationally efficient
    \item Enables much deeper networks
\end{itemize}

\textbf{Problems}:
\begin{itemize}
    \item \textbf{Dying ReLU}: Neurons with negative input have zero gradient, never recover
    \item \textbf{Not zero-centered}: All positive outputs
\end{itemize}

\subsubsection{Leaky ReLU and PReLU}

\begin{align}
\text{LeakyReLU}(x) &= \max(\alpha x, x), \quad \alpha = 0.01 \text{ (typically)} \\
\text{PReLU}(x) &= \max(\alpha x, x), \quad \alpha \text{ learned}
\end{align}

\textbf{Fixes dying ReLU}: Small gradient for negative inputs.

\subsubsection{ELU (Exponential Linear Unit)}

\begin{equation}
\text{ELU}(x) = \begin{cases} x & x > 0 \\ \alpha(e^x - 1) & x \leq 0 \end{cases}
\end{equation}

\textbf{Properties}:
\begin{itemize}
    \item Smooth at $x = 0$
    \item Mean activations closer to zero
    \item Saturates for negative values (noise robustness)
\end{itemize}

\subsection{Modern Activations}

\subsubsection{GELU (Gaussian Error Linear Unit)}

\begin{mathresult}{GELU}
\begin{equation}
\text{GELU}(x) = x \cdot \Phi(x) = x \cdot P(Z \leq x), \quad Z \sim \mathcal{N}(0, 1)
\end{equation}

Approximation:
\begin{equation}
\text{GELU}(x) \approx 0.5x\left(1 + \tanh\left(\sqrt{2/\pi}(x + 0.044715x^3)\right)\right)
\end{equation}
\end{mathresult}

\textbf{Intuition}: Weight input by its percentile under Gaussian distribution.

\textbf{Used in}: BERT, GPT-2, GPT-3, ViT

\subsubsection{SiLU / Swish}

\begin{equation}
\text{SiLU}(x) = x \cdot \sigma(x) = \frac{x}{1 + e^{-x}}
\end{equation}

\textbf{Properties}:
\begin{itemize}
    \item Self-gated: Uses input to gate itself
    \item Smooth, non-monotonic
    \item Slight negative region (regularization effect)
\end{itemize}

\textbf{Used in}: EfficientNet, LLaMA (in SwiGLU)

\subsubsection{SwiGLU (Swish-Gated Linear Unit)}

\begin{mathresult}{SwiGLU}
For FFN in Transformers:
\begin{equation}
\text{SwiGLU}(\vx, \mW, \mV, \mW_2) = (\text{Swish}(\vx\mW) \odot \vx\mV) \mW_2
\end{equation}
\end{mathresult}

\textbf{Key insight}: Multiplicative gating provides richer interactions.

\textbf{Used in}: LLaMA, PaLM, modern LLMs

\textbf{Note}: Requires 3 weight matrices instead of 2, so FFN expansion factor reduced (8/3 instead of 4).

\subsection{Activation Selection Guide}

\begin{table}[H]
\centering
\caption{Activation function selection}
\begin{tabularx}{\textwidth}{lXX}
\toprule
\textbf{Architecture} & \textbf{Recommended} & \textbf{Alternative} \\
\midrule
CNN (general) & ReLU & LeakyReLU \\
BERT-style encoder & GELU & ReLU \\
GPT-style decoder & GELU & SwiGLU \\
Modern LLMs & SwiGLU & GELU \\
Output (classification) & Softmax & Sigmoid (multi-label) \\
Output (regression) & None (linear) & ReLU (non-negative) \\
\bottomrule
\end{tabularx}
\end{table}

%==============================================================================
\section{Normalization Techniques}
\label{sec:normalization}
%==============================================================================

\subsection{Why Normalization?}

\textbf{Internal covariate shift}: Distribution of layer inputs changes during training, making optimization harder.

\textbf{Normalization benefits}:
\begin{itemize}
    \item Stabilizes training
    \item Enables higher learning rates
    \item Reduces sensitivity to initialization
    \item Acts as regularization
\end{itemize}

\subsection{Batch Normalization}

\begin{mathresult}{Batch Normalization}
For a mini-batch $\mathcal{B} = \{x_1, \ldots, x_m\}$:
\begin{align}
\mu_\mathcal{B} &= \frac{1}{m}\sum_{i=1}^m x_i \\
\sigma_\mathcal{B}^2 &= \frac{1}{m}\sum_{i=1}^m (x_i - \mu_\mathcal{B})^2 \\
\hat{x}_i &= \frac{x_i - \mu_\mathcal{B}}{\sqrt{\sigma_\mathcal{B}^2 + \epsilon}} \\
y_i &= \gamma \hat{x}_i + \beta
\end{align}
where $\gamma, \beta$ are learned parameters.
\end{mathresult}

\textbf{Training vs. Inference}:
\begin{itemize}
    \item Training: Use batch statistics
    \item Inference: Use running mean/variance from training
\end{itemize}

\textbf{Limitations}:
\begin{itemize}
    \item Requires sufficiently large batch sizes
    \item Different behavior train vs. test
    \item Problematic for sequence models (variable length)
    \item Not used in Transformers
\end{itemize}

\subsection{Layer Normalization}

\begin{mathresult}{Layer Normalization}
For each sample, normalize across features:
\begin{align}
\mu &= \frac{1}{d}\sum_{i=1}^d x_i \\
\sigma^2 &= \frac{1}{d}\sum_{i=1}^d (x_i - \mu)^2 \\
\hat{x}_i &= \frac{x_i - \mu}{\sqrt{\sigma^2 + \epsilon}} \\
y_i &= \gamma_i \hat{x}_i + \beta_i
\end{align}
\end{mathresult}

\textbf{Key difference from BatchNorm}: Normalizes across features (within sample), not across batch.

\textbf{Advantages}:
\begin{itemize}
    \item Works with any batch size (even 1)
    \item Same behavior train and test
    \item Natural for sequence models
\end{itemize}

\textbf{Used in}: Transformers (BERT, GPT, etc.)

\subsection{RMSNorm}

\begin{mathresult}{RMSNorm}
\begin{equation}
\text{RMSNorm}(\vx) = \frac{\vx}{\text{RMS}(\vx)} \odot \bm{\gamma}, \quad \text{RMS}(\vx) = \sqrt{\frac{1}{d}\sum_{i=1}^d x_i^2}
\end{equation}
\end{mathresult}

\textbf{Simplification of LayerNorm}:
\begin{itemize}
    \item No mean centering (only scale normalization)
    \item No bias term $\beta$
    \item Empirically similar performance
    \item 10-15\% faster
\end{itemize}

\textbf{Used in}: LLaMA, Mistral, most modern LLMs

\subsection{Other Normalization Variants}

\subsubsection{Instance Normalization}

Normalize each channel of each sample independently.
\begin{equation}
\text{IN}(x_{nchw}) = \frac{x_{nchw} - \mu_{nc}}{\sigma_{nc}}
\end{equation}

\textbf{Used in}: Style transfer, image generation

\subsubsection{Group Normalization}

Divide channels into groups, normalize within groups.

\textbf{Used in}: When batch size is very small (object detection)

\subsubsection{Weight Normalization}

Reparameterize weights: $\vw = g \frac{\vv}{\|\vv\|}$

\textbf{Used in}: Some generative models

\subsection{Pre-Norm vs. Post-Norm}

\begin{figure}[H]
\centering
\begin{tikzpicture}[
    node distance=0.5cm,
    box/.style={rectangle, draw, minimum width=2cm, minimum height=0.6cm, rounded corners, font=\small},
    norm/.style={box, fill=lightyellow},
    sublayer/.style={box, fill=lightblue}
]
    % Post-norm
    \begin{scope}[shift={(0,0)}]
        \node[sublayer] (sub1) {SubLayer};
        \node[above=0.3cm of sub1] (add1) {$+$};
        \node[norm, above=0.3cm of add1] (norm1) {Norm};
        \node[above=0.5cm of norm1] (out1) {};
        \node[below=0.5cm of sub1] (in1) {};
        
        \draw[->] (in1) -- (sub1);
        \draw[->] (sub1) -- (add1);
        \draw[->] (add1) -- (norm1);
        \draw[->] (norm1) -- (out1);
        \draw[->] (in1.north) -- ++(-0.8,0) |- (add1.west);
        
        \node[below=1cm of sub1, font=\bfseries] {Post-Norm};
        \node[below=1.4cm of sub1, font=\small] {(Original Transformer)};
    \end{scope}
    
    % Pre-norm
    \begin{scope}[shift={(5,0)}]
        \node[norm] (norm2) {Norm};
        \node[sublayer, above=0.3cm of norm2] (sub2) {SubLayer};
        \node[above=0.3cm of sub2] (add2) {$+$};
        \node[above=0.5cm of add2] (out2) {};
        \node[below=0.5cm of norm2] (in2) {};
        
        \draw[->] (in2) -- (norm2);
        \draw[->] (norm2) -- (sub2);
        \draw[->] (sub2) -- (add2);
        \draw[->] (add2) -- (out2);
        \draw[->] (in2.north) -- ++(-0.8,0) |- (add2.west);
        
        \node[below=1cm of norm2, font=\bfseries] {Pre-Norm};
        \node[below=1.4cm of norm2, font=\small] {(Modern Standard)};
    \end{scope}
\end{tikzpicture}
\caption{Pre-norm vs. post-norm placement in Transformer blocks.}
\end{figure}

\textbf{Pre-norm advantages}:
\begin{itemize}
    \item More stable training (gradient flows through residual directly)
    \item No learning rate warmup needed
    \item Standard in modern LLMs
\end{itemize}

\subsection{Normalization Selection Guide}

\begin{table}[H]
\centering
\caption{Normalization selection by architecture}
\begin{tabularx}{\textwidth}{lXX}
\toprule
\textbf{Architecture} & \textbf{Recommended} & \textbf{Why} \\
\midrule
CNN (large batch) & BatchNorm & Proven, effective \\
CNN (small batch) & GroupNorm & Batch-independent \\
Transformer & LayerNorm & Sequence-friendly \\
Modern LLM & RMSNorm & Faster, similar quality \\
Style transfer & InstanceNorm & Per-image normalization \\
GAN discriminator & SpectralNorm & Stabilizes training \\
\bottomrule
\end{tabularx}
\end{table}

\begin{warningbox}
BatchNorm in multi-GPU training requires explicit synchronization (SyncBatchNorm). Without it, each GPU computes statistics from its local batch, which can be too small and lead to noisy estimates. This is a common source of subtle training bugs when scaling to multiple GPUs---loss looks fine but model quality degrades.
\end{warningbox}

\begin{interviewq}
\textbf{Q: Why don't Transformers use BatchNorm?}

\textbf{What the interviewer is testing:} Understanding of normalization mechanics and why architectural choices are made.

\textbf{A:} Several reasons:

\begin{enumerate}
    \item \textbf{Variable sequence lengths}: BatchNorm statistics would be dominated by padding positions

    \item \textbf{Small effective batch}: With long sequences, effective batch might be small (sequence length $\times$ batch size tokens, but each position is different)

    \item \textbf{Train/test discrepancy}: Running statistics problematic for autoregressive generation

    \item \textbf{LayerNorm works well}: Normalizing across features (within token) aligns with attention's token-wise processing
\end{enumerate}

LayerNorm normalizes each token independently, which matches how attention processes sequences.

\textbf{Common follow-ups:} ``Why RMSNorm over LayerNorm in modern LLMs?'' (Drops mean-centering, 10-15\% faster, similar quality.) ``Where is BatchNorm still preferred?'' (CNNs with large batches, especially for vision tasks.)
\end{interviewq}

\begin{interviewq}
\textbf{Q: Your model's training loss oscillates wildly and doesn't converge. What activation/normalization issues might cause this?}

\textbf{What the interviewer is testing:} Debugging intuition, connecting architectural choices to training dynamics.

\textbf{A:} I'd check these activation/normalization issues in order:

\begin{enumerate}
    \item \textbf{No normalization or wrong placement}: Pre-norm is much more stable than post-norm for deep networks. Adding LayerNorm before attention/FFN (pre-norm) often fixes instability.

    \item \textbf{Dying ReLU}: If using ReLU, check what fraction of activations are zero. If $>$ 50\% of neurons are dead, switch to GELU or LeakyReLU.

    \item \textbf{BatchNorm with small batches}: If batch size per GPU is small ($<$ 16), BatchNorm statistics are noisy. Switch to GroupNorm or LayerNorm.

    \item \textbf{Numerical overflow in softmax}: Large attention logits before softmax can cause NaN. This is why we scale by $1/\sqrt{d_k}$. Check if attention weights contain NaN/Inf.

    \item \textbf{Learning rate too high}: Normalization enables higher learning rates, but there's still a limit. Try halving the learning rate first.
\end{enumerate}

\textbf{Common follow-ups:} ``How do you diagnose dying ReLU in practice?'' (Histogram of activations; check gradient norms per layer.) ``What's the relationship between normalization and learning rate?'' (Normalization makes the loss landscape smoother, enabling larger learning rates.)
\end{interviewq}